\documentclass[11pt]{article}
\usepackage{amsmath,amsfonts,mathabx}
\usepackage{graphicx}
\usepackage{hyperref}

% PREAMBLE
\newcommand{\RR}{\ensuremath{\mathbb{R}}}
\newcommand{\RRn}[1]{\mathbb{R}^{#1}}
\DeclareMathOperator{\function}{function}

% Code listing
\usepackage[listings]{tcolorbox}
\usepackage[normalem]{ulem}
\usepackage{xcolor}


\begin{document}


\begin{center}
{\textbf{\LARGE Chapitre 9: R\'egression non-lin\'eaire}}
\end{center}

\noindent Consid\'erons le jeu de donn\'ees \textsf{Auto} de la library \texttt{ISLR} du logiciel \textrm{R}. Ce jeu de donn\'ees contient 392 observations. La variable r\'eponse est \texttt{mpg}, la distance moyenne en miles parcourue par la voiture avec un gallon de carburant. On s'int\'eresse \`a une seule variable explicative, soit le poids de la voiture \textrm{weight}.

\noindent On commence par tracer le nuage de points $\{(X_i,Y_i), i=1,\cdots,392\}$, o\`u $X_i$ et $Y_i$ sont, respectivement, le poids et la distance moyenne en miles parcourue avec un gallon de carburant de la voiture $i$. 

\begin{verbatim}
> if (!require("ISLR")) install.packages("ISLR")
> library("ISLR")#### Installer la librairie ISLR
> 
> data(Auto)
> X <- Auto$weight
> Y <- Auto$mpg
> 
> plot(X,Y,xlab="Poids de la voiture",
+      ylab="Distance parcourue avec un gallon de carburant")
\end{verbatim}

\noindent Et on obtient 
\begin{center}
\includegraphics[scale=0.3]{Graphique1.pdf}
\end{center}


\noindent Ce graphique montre que la relation entre $X$ et $Y$ n'est clairement pas \sout{clairement pas }lin\'eaire. Dans le chapitre pr\'ec\'edent, on a consid\'er\'e deux \sout{approaches}\textcolor{red}{approches} pour \'etudier la relation entre $X$ et $Y$; soient l'estimateur des $k$ plus proches voisins et l'estimateur par noyau. Ces approches nous ont permis d'estimer $f(x)=E[Y|X=x]$. Dans ce chapitre, on pr\'esente des approches al\sout{e}\textcolor{red}{}ternatives pour estimer $f(x)$.

\section{R\'egression polynomiale}

\noindent L'id\'ee de la r\'egression polynomiale est d'approximer $f(x)$ par un polyn\^ome de degr\'e $p$.  Ainsi, on pose 
$$f(x) = \beta_0 + \beta_1 x + \cdots + \beta_p x\textcolor{red}{^p}$$
et on obtient le mod\`ele lin\'eaire 
$$Y_i = \beta_0 + \beta_1 X_i + \cdots + \beta_p X_i^p + \epsilon_i, \, \, \, i=1,\cdots,n.$$
\noindent On peut faire l'hypoth\`ese que $\epsilon_i \sim {\cal N}(0,\sigma^2)$ ou que simplement $E[\epsilon_i]=0$ et $Var[\epsilon_i]=\sigma^2>0$.

\noindent Dans tous les cas, on estime $\beta=(\beta_0,\beta_1,\cdots,\beta_p)^\top$ par 
$$\hat \beta = (X^\top X)^{-1} X^\top Y,$$
o\`u
$$X =\begin{pmatrix} 1 & X_1 & \cdots & X_1^p \\
                     1 & X_2 & \cdots & X_2^p \\
                     \vdots & \vdots & \ddots & \vdots \\
                     1 & X_n & \cdots & X_n^p \\
\end{pmatrix}.$$ 

\noindent Le vecteur $\beta$ de coefficients de regr\'ession contient $p+1$ \'el\'ements. On dit qu'on a $p+1$ degr\'es de libert\'e.

\noindent Pour s\'electioner la valeur la plus ad\'equate pour $p$, on peut se fier \`a des crit\`eres comme l'AIC ou le BIC ou encore effectuer des tests d'hypoth\`eses avec des statistiques $F$. Dans ce chapitre, on adopte l'esprit du chapitre pr\'ec\'edent et on choisi\textcolor{red}{t} la valeur de $p$ qui minimise le taux d'erreur test. On peut estimer le taux d'erreur test avec la proc\'edure qui exclut une observation du jeu de donn\'ees entrainement \`a la fois. On a vu au Chapitre 7 que cette m\'ethode nous conduit \`a une expression explicite pour le crit\`ere de la validation crois\'ee.

\noindent Une fois qu'on choisi\textcolor{red}{t} la valeur de $p$, on a alors
$$\hat f(x) = \hat\beta_0 + \hat\beta_1 x + \cdots + \sout{\beta_p}\textcolor{red}{\hat\beta_p} x^p = x_* \hat\beta;$$
\noindent o\`u $x_* = \begin{pmatrix} 1 & x & x^2 & \cdots & \sout{x_p}\textcolor{red}{x^p} \end{pmatrix}$.
\noindent On en d\'eduit que $Var[\hat f(x)]= x_* Var[\hat\beta] x_*^\top$. Cette expression nous permet d'obtenir des bandes de confiance autour de $\hat f(x)$.

\begin{verbatim}
> ############################
> ## Regression polynomiale ##
> ############################
> 
> modele.polynomial <- lm(mpg~ poly(weight,2,raw=TRUE),data=Auto)
> 
> axe.poids <- seq(min(Auto$weight),max(Auto$weight),1)
> prediction.polynomial <- predict(modele.polynomial,
>                                  newdata=list(weight=axe.poids),se=TRUE)
> 
> plot(mpg~weight,data=Auto)
> lines(axe.poids,prediction.polynomial$fit,col="red")
> lines(axe.poids,
> prediction.polynomial$fit-1.96*prediction.polynomial$se.fit,lty=2,col="blue")
> lines(axe.poids,
> prediction.polynomial$fit+1.96*prediction.polynomial$se.fit,lty=2,col="blue")
\end{verbatim}

\noindent Et on obtient 
\begin{center}
\includegraphics[scale=0.3]{Graphique2.pdf}
\end{center}



\noindent La regr\'ession polynomiale peut \^etre vue comme une regr\'ession sur des transformations connues de la variable explicative: 
$$Y_i = \beta_0 b_0(X_i)+ \beta_1 b_1(X_i) + \cdots + \beta_p b_p(X_i) + \epsilon_i, \, \, \, i=1,\cdots,n,$$

\noindent avec $b_k(x)=x^k$ pour $k=0,1,\cdots,p$. Cette approche peut \^etre g\'en\'eralis\'ee de plusieurs fa\c{c}ons.

\section{R\'egression avec des fonctions constantes par mo\sout{c}\textcolor{red}{r}ceaux}

Posons $[A,B]$ le domaine de d\'efinition de la variable explicative $X$. Dans la pratique, on peut par exemple, poser $A=\mbox{min}\{X_1,\cdots,X_n\}-\varepsilon$ et $B=\mbox{\sout{min}}\textcolor{red}{max}\{X_1,\cdots,X_n\}+\varepsilon$ o\`u $\varepsilon$ est un tr\`es petit r\'eel positif. 

\noindent Consid\'erons une partition $A = c_0 < c_1 < \cdots < c_K = B$ de l'intervalle $[A,B]$. Pour $k=1,\cdots,K$, on d\'efi\sout{tion}\textcolor{red}{nit} la fonction indicatrice $b_k(x)$ par:
$$b_k(x) = \begin{cases} 1 & \mbox{ si } x \in [c_{k-1},c_k] \\ 0 & \mbox{ sinon }\end{cases}$$

\noindent Cette approche approxime $f(x)$ par une fonction constante par escalier: $f(x)=\beta_k$ si $x \in [c_{k-1},c_k]$. Autrement dit, on a
$$f(x) = \beta_1 b_1(x) + \beta_2 b_2(x) + \cdots + \beta_K b_K(x)$$
\noindent et le mod\`ele devient 
$$Y_i = \beta_1 b_1(X_i) + \beta_2 b_2(X_i) + \cdots + \beta_K b_K(X_i) + \epsilon_i$$
\noindent La variable explicative $X$ a donc \'et\'e discr\'etis\'ee et on a un mod\`ele d'ANOVA \`a un facteur fixe \`a $K$ modalit\'es. On montre que l'estimateur des moindres carr\'es de $\beta_k$ est
$$\hat\beta_k = \frac{\sum\limits_{i=1}^n b_k(X_i)Y_i}{n_k},$$
\noindent o\`u $n_k=\sum\limits_{i=1}^n b_k(X_i)$ est le nombre d'observations $X_i$ qui appartiennent \`a $[c_{k-1},c_k]$. Finalement, on a $Var[\hat\beta_k]=\sigma^2/n_k$. 

\begin{verbatim}
> #######################################
> ## Regression constante par morceaux ##
> #######################################
> 
> modele.par.morceaux <- lm(mpg~ cut(weight,4),data=Auto)
> prediction.par.morceaux <- predict(modele.par.morceaux,
>                                    newdata=list(weight=axe.poids),se=TRUE)
> 
> plot(mpg~weight,data=Auto)
> lines(axe.poids,prediction.par.morceaux$fit,col="red")
> lines(axe.poids,
> prediction.par.morceaux$fit-1.96*prediction.par.morceaux$se.fit,lty=2,col="blue")
> lines(axe.poids,
> prediction.par.morceaux$fit+1.96*prediction.par.morceaux$se.fit,lty=2,col="blue")
\end{verbatim}

\noindent Et on obtient 
\begin{center}
\includegraphics[scale=0.3]{Graphique3.pdf}
\end{center}


\section{R\'egression avec des splines}

Les splines sont des polyn\^omes d\'efinis par intervalles. On commence par un cas simple avec $K=2$ et $p=3$. On divise le domaine $[A,B]$ en deux intervalles : $[c_0,c_1)$ et $[c_1,c_2)$. On approxime la fonction $f$ par 
$$
f(x) = \begin{cases} \beta_0 + \beta_1 x + \beta_2 x^2 + \beta_3 x^3 & \mbox{ si } x \in [c_0,c_1), \\ 
                     \gamma_0 + \gamma_1 x + \gamma_2 x^2 + \gamma_3 x^3 & \mbox{ si } x \in [c_1,c_2). \end{cases}
$$

\noindent On a donc $4+4=8$ degr\'es de libert\'e. En plus, on impose que $f$ soit continue et deux fois d\'erivable au point $c_1$. Ce qui nous donne 3 contraintes. Finalement, on se retrouve avec 8-3=5 degr\'es de libert\'e.

\noindent G\'en\'eralisons ce qu'on vient de faire \`a $K$ et $p$ arbitraires. On a $K$ polyn\^omes de degr\'e $p$, ce qui nous donne $K\times (p+1)$ param\`etres. Aux $K-1$ points int\'erieurs, on impose que $f$ soit continue et $p-1$ fois d\'erivable, ce qui nous donne $p\times (K-1)$ contraintes. Finalement, on a $K(p+1)-(K-1)p=K+p$ degr\'es de libert\'e.

\noindent Pour un $p$ et un $K$ donn\'es, un choix math\'ematiquement convivial pour construire une base de $K+p$ splines est donn\'e par: $b_1(x)=1, b_2(x)=x, \cdots, b_{p+1}(x)=x^p, b_{p+2}(x)=h(x,c_1), b_{p+3}=h(x,c_2), \cdots, b_{p+K}=h(x,c_{K-1}),$

\noindent o\`u $$h(x,c)=\begin{cases} (x-c)^p & \mbox{ si } x \geq c  \\ 0 \mbox{ sinon.} \end{cases}$$

\noindent Il existe une infinit\'e de choix pour construire une base de $K+p$ splines. En effet, n'importe quelles $K+p$ combinaisons lin\'eaires ind\'ependantes de la base qu'on vient de d\'efinir est encore une base de $K+p$ splines.

\noindent On d\'efinit alors le mod\`ele 
$$Y_i  = \beta_1 b_1(X_i) + \beta_2 b_2(X_i) + \cdots + \beta_{K+p} b_{K+p}(X_i) + \epsilon_i, \, \, \, i=1,\cdots,n.$$ 

\noindent L'estimateur des moindres carr\'es de $\beta=(\sout{beta_1}\textcolor{red}{\beta_1},\cdots,\beta_{K+p})^\top$ est $(X^\top X)^{-1} X^\top Y,$ o\`u 
$$X =
\begin{pmatrix}
b_1(X_1) & b_2(X_1) & \cdots & b_{K+p}(X_1) \\
b_1(X_2) & b_2(X_2) & \cdots & b_{K+p}(X_2) \\
\vdots & \vdots & \ddots & \vdots \\
b_1(X_n) & b_2(X_n) & \cdots & b_{K+p}(X_n) 
\end{pmatrix}.
$$

\noindent On obtient finalement $\hat f(x) = \hat\beta_1 b_1(x) + \hat\beta_2 b_2(x) + \cdots + \hat\beta_{K+p} b_{K+p}(x)$.

\begin{verbatim}
######################
## Splines cubiques ##
######################

if (!require("splines")) install.packages("splines")
library("splines")	#### Installer la librairie splines

modele.splines <- lm(mpg~bs(weight,knots=c(2200,2800,3600)),data=Auto)
prediction.splines <- predict(modele.splines,
                              newdata=list(weight=axe.poids),se=TRUE)
plot(mpg~weight,data=Auto)
lines(axe.poids,prediction.splines$fit,col="red")
lines(axe.poids,
prediction.splines$fit+1.96*prediction.splines$se.fit,lty=2,col="blue")
lines(axe.poids,
prediction.splines$fit-1.96*prediction.splines$se.fit,lty=2,col="blue")
\end{verbatim}

\noindent Et on obtient 
\begin{center}
\includegraphics[scale=0.3]{Graphique4.pdf}
\end{center}

\noindent Cet estimateur peut \^etre instable aux extr\'emit\'es du domaine de $X$. Pour rem\'edier \`a ce probl\`eme, on consid\`ere les splines naturelles. Celles-ci sont d\'efinies exactement comme les splines, \`a l'exception qu'on consid\`ere seulement des fonctions lin\'eaires sur les intervalles $[c_0,c_1)$ et $[c_{K-1},c_K)$.

\begin{verbatim}

#################################
## Splines cubiques naturelles ##
#################################

modele.sn <- lm(mpg~ns(weight,df=4),data=Auto)
prediction.sn <- predict(modele.sn,newdata=list(weight=axe.poids),se=TRUE)
plot(mpg~weight,data=Auto)
lines(axe.poids,prediction.sn$fit,col="blue")
lines(axe.poids,
prediction.sn$fit+1.96*prediction.sn$se.fit,lty=2,col="blue")
lines(axe.poids,
prediction.sn$fit-1.96*prediction.sn$se.fit,lty=2,col="blue")
\end{verbatim}

\noindent Et on obtient 
\begin{center}
\includegraphics[scale=0.3]{Graphique5.pdf}
\end{center}

\noindent Un autre d\'efau\sout{l}t des splines est qu'elles peuvent avoir tendance \`a aller chercher le maximum de points. Pour rem\'edier \`a ce probl\`eme, on impose une p\'enalit\'e pour les splines trop variables. Ainsi, on estime $\beta$ en minimisant
$$\sum_{i=1}^n \left\{Y_i - [\beta_1 b_1(X_i) + \beta_2 b_2(X_i) + \cdots + \beta_{K+p} b_{K+p}(X_i)]\right\}^2 + \lambda \int_A^B f''(t)^2 dt.$$

\noindent On montre que l'estimateur qui en r\'esulte est donn\'e par
$$\hat \beta_\lambda = (X^\top X + \lambda M)^{-1} X^\top Y.$$

\noindent Ici, $M$ est une matrice carr\'e\textcolor{red}{e} de dimensi\sout{oo}\textcolor{red}{o}n $K+p$ et l'\'el\'ement \`a la position $(l,q)$ de la matrice $M$ est \'egal \`a $\int_A^B b''_l(t) b''_q(t) dt$.

\noindent Le param\`etre $\lambda$ est un param\`etre qui contr\^ole la r\'egularit\'e de $\hat f$. Dans la pratique, on estime ce param\`etre par validation crois\'ee.

\begin{verbatim}
> #######################
> ## Splines penalites ##
> #######################
> 
> modele.sp.df.16 <- smooth.spline(Auto$weight,Auto$mpg,df=16)
> 
> plot(mpg~weight,data=Auto)
> lines(modele.sp.df.16,col="green")
> 
> modele.sp.cv <- smooth.spline(Auto$weight,Auto$mpg,cv=TRUE)
Message d'avis :
Dans smooth.spline(Auto$weight, Auto$mpg, cv = TRUE) :
  la validation croisée avec des valeurs de 'x' non uniques paraît douteuse
> modele.sp.cv$df
[1] 10.35561
> lines(modele.sp.cv,col="purple")
\end{verbatim}

\noindent Et on obtient 
\begin{center}
\includegraphics[scale=0.3]{Graphique6.pdf}
\end{center}

\end{document}